\documentclass[a4paper,12pt]{letter}
\usepackage{geometry}
\usepackage{graphicx}
\geometry{a4paper, margin=1in}

\begin{document}

% Letterhead 
\noindent
[Your Department/Institute Name] \\
[Your Institution Name]

\vspace{0.5em}
\noindent
[City, Country], \today

\vspace{2em}

% Add institution logo if available - right aligned
% \begin{flushright}
% \includegraphics[width=0.3\textwidth]{institution_logo.png}
% \end{flushright}

\vspace{1em}

\noindent
To: The Editors, \\
Nature Communications Physics

\vspace{1em}

\noindent
\textbf{Re: Submission of Manuscript titled "Guided rewiring of social networks reduces polarization and accelerates collective action"}

\vspace{1em}

\noindent
Dear Editors,

\vspace{0.5em}

On behalf of my co-authors, I am pleased to submit our article "Guided rewiring of social networks reduces polarization and accelerates collective action" for your consideration for publication in Nature Communications Physics. In this paper, we address the pressing question of how to rapidly depolarize social networks and encourage cooperative consensus on large-scale collective action challenges.

We confirm that this work is original and has not been published elsewhere, nor is it under consideration for publication elsewhere.

Much recent scientific work and public discourse has focused on algorithmic influence on social interactions and community formation in online social media, and more broadly in digital common spaces. Evidence shows that contemporary algorithmic design, as utilized by large online platforms such as X (formerly Twitter), Facebook, and TikTok, may work to polarize participants in digital common spaces and thus hinder cooperative consensus formation.

Previous empirical and modeling work has investigated the effect of commercially deployed link-recommendation algorithms on social network inequalities and, separately, the effect of idealized link-recommendations on opinion dynamics. A gap exists, however, in the combination of these commercially deployed algorithms with classical opinion dynamics models. Additionally, the existence of the backfire effect—where social interactions increase, rather than reduce, the gap between opinions—remains disputed.

We formulate two novel rewiring algorithms based on principles identified in the literature to reduce polarization and compare these against two commonly used rewiring algorithms employed by platforms such as Twitter (now X) across a number of real and idealized network topologies. We demonstrate that one of the commercial algorithms, Who To Follow, generated and maintained polarization across empirical and synthetic networks. Furthermore, it prohibited cooperative consensus formation even under the influence of a cooperative field effect.

We believe that the timeliness of this topic and the originality of our approach will appeal to the community working within the scope of Nature Communications Physics. Our work contributes to the literature in that it (1) combines an opinion dynamics model with commercial-use topological rewiring algorithms derived from Who To Follow and node2vec, (2) tests this framework on real-world network structures, (3) provides algorithms that, per our results, are more effective than existing commercial algorithms at reducing polarization, and (4) includes both directed and undirected networks in our analysis. Our results have the potential to inform strategies for leveraging social media as powerful tools for accelerating collective action on urgent global issues such as climate change.

Our work addresses the intersection of network theory, complex systems, statistical physics, and computational social science, providing quantitative insights into how to resolve the polarization that currently impedes effective political and prosocial collective action on large-scale challenges. We believe this aligns well with the scope and readership of Nature Communications Physics.

Please do not hesitate to contact me if you have any questions or if you need any further information.

\vspace{1em}

\noindent
Yours sincerely,

\vspace{2em}

\noindent
Jordan Everall \\
[Your Position/Title] \\
[Your Department/Institute Name] \\
[Your Institution Name] \\
[City, Country] \\
Email: [your.email@institution.edu]

\end{document}
